\section{Discussion and Conclusion}
\label{sec:discussion}
The three studies distinguish practical benefit from evidence about transfer. Phase~I shows that an analogy condition can improve average accuracy without outperforming an explanation control, and can repair some answers while damaging others. The benefit is real under that comparison, but its interpretation depends on what the baseline controls. Textual mention adds a behavioral observation without resolving whether the relevant relation was used correctly.

Phase~II provides a more targeted test. The source answer helps several models more when attached to a structurally corresponding source than to a reassigned source from the same corpus. This supports structure-dependent use of source solutions under the tested construction and prompting protocol. It also distinguishes three behaviors: index following, benefits from corresponding source content without a selective answer effect, and selective source-answer benefits. A single accuracy ranking would obscure these differences.

The framework connects these outcome distinctions to inspectable operations. Its central contribution is the combination of theoretical stage distinctions, explicit information contracts, and continuation tests after intermediate corrections. Cognitive theory motivates evaluating source access, relational alignment, abstraction, and target application separately. The workflow makes these distinctions executable without asserting that its generated intermediates expose the model's internal computation.

The repairs reveal a useful diagnostic distinction: \emph{where a path first fails} differs from \emph{which capabilities remain inadequate after repair}. GPQA initially appears to have a retrieval-only problem, yet most of its failed paths remain unsuccessful after useful upstream information is supplied. For murder mysteries, adaptation is prominent both initially and after upstream repair. Thus, better sources alone are insufficient; a model must preserve the useful structure and apply it under the target's constraints.

Future work can compare this information flow with alternative stage interfaces, expand source validation, and test robustness across model families and domains. Matched malformed-stage controls would further separate the value of a specific correction from generic benefits of additional guidance. The present evidence supports an evaluation principle: measure helpful and harmful outcome changes, test whether source-answer benefits depend on correspondence, and inspect the stages needed to turn a source into a target solution. This gives a more precise account of analogical reasoning than accuracy or mapping quality alone.
